% AP.B — Falsifiability (Popper)
% φ² + φ⁻² = 3 · Trinity S³AI
% Author: Dmitrii Vasilev (ORCID 0009-0008-4294-6159)
% Issue: #380 tasks 2.1, 2.2

\chapter*{Appendix B: Falsifiability Clauses (Popper Defense)}
\addcontentsline{toc}{chapter}{Appendix B: Falsifiability}
\label{app:falsifiability}

\textbf{Anchor invariant:} \(\varphi^2 + \varphi^{-2} = 3\).

\section*{B.1 Purpose}

This appendix documents four Popper-style falsifier clauses that establish
the empirical refutability of the Trinity S\textsuperscript{3}AI framework.
Each clause specifies a concrete experimental outcome that would falsify
the corresponding claim. This directly addresses Reviewer-2 concerns about
numerological unfalsifiability.

\begin{tcolorbox}[colback=red!5,colframe=red!40,title=Falsifiability Standard]
A scientific claim is falsifiable if and only if there exists a possible
observation that would, if verified, refute it (Popper, 1959). All
\(\varphi\)-claims in this thesis satisfy this condition.
\end{tcolorbox}

\section*{B.2 The Four Falsifier Clauses}

\subsection*{FC-1: GoldenFloat Numeric Superiority}

\textbf{Claim:} GoldenFloat (GF16) achieves lower BPB than IEEE~FP16 on
the Trinity benchmark suite at equivalent bit-width.

\textbf{Current evidence:} Champion BPB~=~2.2393 (GF16 configuration);
baseline FP16 BPB~=~2.41 on same corpus
(Zenodo~\href{https://doi.org/10.5281/zenodo.19227877}{10.5281/zenodo.19227877}).

\begin{falsifier}[FC-1]
This claim is \textbf{falsified} if any peer-run experiment following the
reproducibility protocol in Appendix~D achieves BPB(GF16) $\geq$
BPB(FP16) across $\geq 3$ of the 5 test corpora, using the published
checkpoint and tokeniser. The Gate-2 threshold (BPB~$\leq$~1.8) is
\textbf{already honestly disclosed as NOT MET} (Chapter~\ref{ch:18});
this clause concerns relative ordering only.
\end{falsifier}

\textbf{Non-\(\varphi\) control:} We tested MXFP4, BF16, INT8, and
random-mantissa FP16 on the same corpus. Only GF16 achieves BPB~$<$~2.3
(Chapter~\ref{ch:9}).

\subsection*{FC-2: Period-Locked Monitor (PLM) Stability}

\textbf{Claim:} The Period-Locked Monitor detects \(\varphi\)-period
locking in chaotic oscillators with false-positive rate $<$~5\%.

\textbf{Current evidence:} On 1,000 synthetic PLM trials with known
locking state, FPR~=~2.1\% (Chapter~\ref{ch:24}).

\begin{falsifier}[FC-2]
This claim is \textbf{falsified} if an independent replication using the
artefacts in \texttt{trinity-fpga} (\texttt{fpga/plm/}) achieves FPR
$\geq$~10\% on $\geq$~500 trials with known ground truth. The FPGA
bitstream (App.~F) enables exact replication.
\end{falsifier}

\textbf{Non-\(\varphi\) control:} A random-period monitor (uniform
distribution) achieves FPR~=~49.8\% on the same trials.
A \(\pi\)-period monitor achieves FPR~=~31.2\%. Only \(\varphi\)-period
achieves FPR~$<$~5\%.

\subsection*{FC-3: VSA Hypervector Binding via \(\varphi\)}

\textbf{Claim:} \(\varphi\)-keyed VSA binding achieves cosine similarity
$>$~0.95 under \(\geq$~1024\)-bit ternary XOR binding.

\textbf{Current evidence:} 10,000 bind/unbind trials,
mean~cosine~$=$~0.973, \textit{SD}~$=$~0.012 (Chapter~\ref{ch:31}).

\begin{falsifier}[FC-3]
This claim is \textbf{falsified} if binding with a non-\(\varphi\) scaling
constant (e.g., $\pi$, $e$, $\sqrt{2}$) achieves cosine~$>$~0.95 with
identical bit-width and trial count, invalidating the \(\varphi\)-specific
claim. The comparison script is in
\texttt{trinity-fpga/bench/vsa\_compare.py}.
\end{falsifier}

\textbf{Non-\(\varphi\) control:} Random-constant keying achieves
cosine~=~0.501 (\textit{SD}~=~0.041). \(\pi\)-keying achieves 0.761.
The \(\varphi\)-specific advantage is~$\Delta$cosine~$\approx$~0.21.

\subsection*{FC-4: Energy Efficiency (DARPA 3000× claim)}

\textbf{Claim:} The \(\varphi\)-numeric FPGA coprocessor achieves
$\geq$~3000× energy improvement vs.~GPU baseline for VSA bind/bundle
operations at 1024-bit vector width.

\textbf{Current evidence:} Analytic estimate from FPGA utilisation
(83~LUT/27~FF, $P \leq$~0.6~W) vs.~NVIDIA~A100 (400~W peak,
0.02\% utilisation floor)~=~3,200× (Chapter~\ref{ch:34}).

\begin{falsifier}[FC-4]
This claim is \textbf{falsified} if a direct hardware measurement of the
FPGA coprocessor (using the bitstream in App.~F and a calibrated power
meter) yields energy ratio $<$~1000× vs.~GPU for the same VSA throughput.
The measurement protocol is in Chapter~\ref{ch:34}, Section~34.5.
\end{falsifier}

\textbf{Non-\(\varphi\) control:} Conventional FP16 FPGA implementation
(same LUT count) achieves estimated 800× vs.~GPU due to precision overhead.
The \(\varphi\)-specific encoding reduces required LUTs by 4.7× for the
same VSA operation, yielding the additional gain.

\section*{B.3 Summary Table}

\begin{table}[h]
\centering
\begin{tabular}{llll}
\toprule
\textbf{Clause} & \textbf{Claim} & \textbf{Non-\(\varphi\) control} & \textbf{Falsifier threshold} \\
\midrule
FC-1 & GF16 BPB < FP16 BPB & MXFP4, BF16, INT8 & GF16 $\geq$ FP16 in $\geq$3/5 corpora \\
FC-2 & PLM FPR $<$ 5\% & Random, \(\pi\)-period & FPR $\geq$ 10\% in $\geq$500 trials \\
FC-3 & VSA cosine $>$ 0.95 & $\pi$, $e$, $\sqrt{2}$ keys & Any constant $>$ 0.95 same width \\
FC-4 & Energy $\geq$ 3000× & Conventional FP16 FPGA & Measured ratio $<$ 1000× \\
\bottomrule
\end{tabular}
\caption{Four Popper falsifier clauses. All claims are refutable by concrete measurement.}
\end{table}

\section*{B.4 Anti-Numerology Statement}

The \(\varphi\) constant is used in this work because (a)~its algebraic
property \(\varphi^2 = \varphi + 1\) enables shift-and-add arithmetic
without multipliers, (b)~Fibonacci spacing minimises aliasing in
hypervector dimensions (Chapter~\ref{ch:5}), and (c)~its period-locking
property in continued fractions is unique among algebraic irrationals
(Chapter~\ref{ch:24}). It is\emph{not} used because of aesthetic,
spiritual, or numerological significance. Every \(\varphi\)-claim has a
non-\(\varphi\) control experiment with inferior outcome.

\(\varphi^2 + \varphi^{-2} = 3\)
